% ============================================================
%  Section: System Model and Problem Formulation
%  (可直接粘贴到 Overleaf 主文档中; 需要 amsmath, amssymb 宏包)
% ============================================================

\section{SYSTEM MODEL AND PROBLEM FORMULATION}
\label{sec:system_model}

% ==============================================================
\subsection{Problem Statement}
\label{subsec:problem_statement}

We consider the multi-target cooperative coverage problem in a three-dimensional
environment composed of $N$~homogeneous UAVs
$\mathcal{U} = \{U_1, \dots, U_N\}$,
$M$~mobile ground targets $\mathcal{G} = \{G_1, \dots, G_M\}$, and
$K$~dynamic obstacle zones $\mathcal{T} = \{T_1, \dots, T_K\}$, collectively
referred to as entities
$\mathcal{E} = \{\mathcal{U}, \mathcal{G}, \mathcal{T}\}$.
At the beginning of each episode, the positions of all entities are randomly
initialised within a bounded task region.
All UAVs share identical physical parameters, with a fixed observation
radius~$R_{\mathrm{obs}}$ (maximum sensing range) and a coverage
radius~$R_{\mathrm{cov}}$ (proximity required to cover a target).
The system objective is to \emph{maximise the cumulative target coverage within
a finite time horizon~$T_{\max}$ while ensuring collision avoidance}.

% ==============================================================
\subsection{System Model}
\label{subsec:system_model}

We consider a bounded three-dimensional task region
$\mathcal{W} = [-L/2, L/2]^{2} \times [0, H_{\mathrm{w}}]$, where $L$ denotes
the side length of the square operating area and $H_{\mathrm{w}}$ is the height
of the enclosing rigid boundaries.
The environment consists of $N$~UAVs, $M$~targets, and $K$~obstacle zones.
UAV $U_i$, $i \in \{1, \dots, N\}$, has position
$\mathbf{p}_i^{(t)} = (x_i^{(t)}, y_i^{(t)}, z_i^{(t)})^{\top}$ in the world
frame and velocity
$\mathbf{v}_i^{(t)} = (\dot{x}_i^{(t)}, \dot{y}_i^{(t)},
\dot{z}_i^{(t)})^{\top} \in \mathbb{R}^3$ at time step~$t$.
The speed of every UAV is bounded by $v_{\max}$, i.e.\
$\lVert\mathbf{v}_i^{(t)}\rVert \le v_{\max}$.

\subsubsection{Quadrotor Dynamics}

Each UAV is modelled as a quadrotor with mass~$m$, moment of inertia matrix
$\mathbf{J} \in \mathbb{R}^{3 \times 3}$, body radius~$r_{\mathrm{uav}}$, and
thrust coefficient~$k_f$.
The orientation is described by the Euler angles
$\boldsymbol{\Phi}_i = (\phi_i, \theta_i, \psi_i)^{\top}$, representing roll,
pitch, and yaw, respectively.
The four rotors ($l = 1,\dots,4$) rotate at angular speeds~$\omega_l$ and
produce thrusts $f_l = k_f\,\omega_l^2$.
The translational and rotational dynamics in the world frame are
\begin{align}
  m\,\ddot{\mathbf{p}}_i &=
    \mathbf{R}(\boldsymbol{\Phi}_i)\,
    \begin{pmatrix}0\\0\\\sum_{l=1}^{4} f_l\end{pmatrix}
    - mg\,\mathbf{e}_3,
  \label{eq:trans_dyn} \\
  \mathbf{J}\,\dot{\boldsymbol{\omega}}_i &= \boldsymbol{\tau}_i
    - \boldsymbol{\omega}_i \times \mathbf{J}\,\boldsymbol{\omega}_i,
  \label{eq:rot_dyn}
\end{align}
where $\mathbf{R}(\boldsymbol{\Phi}_i) \in SO(3)$ is the rotation matrix from
the body frame to the world frame, $g$ is the gravitational acceleration,
$\mathbf{e}_3 = (0,0,1)^{\top}$ is the unit vector along the world $z$-axis,
$\boldsymbol{\omega}_i \in \mathbb{R}^{3}$ is the body angular velocity, and
$\boldsymbol{\tau}_i = (\tau_\phi, \tau_\theta, \tau_\psi)^{\top}$ is the
body-frame torque vector derived from the rotor thrusts through the allocation
matrix.

\subsubsection{Cascaded PID Controller}

A cascaded PID controller maps the high-level velocity command to the rotor
speeds through three stages: a position loop, an attitude loop, and a motor
mixer.

\paragraph{Velocity command.}
Given the action
$\mathbf{a}_i = (a_i^{x}, a_i^{y})^{\top} \in [-1,1]^{2}$ output by the
policy, the desired velocity is
\begin{equation}
  \mathbf{v}_i^{\mathrm{cmd}} = v_{\max}\,
  \frac{\mathbf{a}_i}{\lVert\mathbf{a}_i\rVert},
  \label{eq:vel_cmd}
\end{equation}
where $v_{\max}$ is the maximum flight speed.
Only the $xy$-components are commanded; the altitude is held constant at
$z_0 = H_{\mathrm{w}}/2$.
The reference position is updated each control step as
$\mathbf{p}_i^{\mathrm{ref}} = \mathbf{p}_i^{(t)} +
\mathbf{v}_i^{\mathrm{cmd}} \cdot \Delta t$,
where $\Delta t$ is the control period.

\paragraph{Position loop.}
A PID controller computes the desired thrust vector from position and velocity
errors:
\begin{equation}
  \mathbf{F}_i^{\mathrm{des}} =
    \mathbf{K}_p^{\mathrm{pos}}\,\mathbf{e}_{p}
    + \mathbf{K}_i^{\mathrm{pos}}\!\int\!\mathbf{e}_{p}\,\mathrm{d}t
    + \mathbf{K}_d^{\mathrm{pos}}\,\mathbf{e}_{v}
    + mg\,\mathbf{e}_3,
  \label{eq:pid_pos}
\end{equation}
where $\mathbf{e}_{p} = \mathbf{p}_i^{\mathrm{ref}} - \mathbf{p}_i$ is the
position error, $\mathbf{e}_{v} = \mathbf{v}_i^{\mathrm{cmd}} - \mathbf{v}_i$
is the velocity error, and
$\mathbf{K}_p^{\mathrm{pos}}, \mathbf{K}_i^{\mathrm{pos}},
\mathbf{K}_d^{\mathrm{pos}} \in \mathbb{R}^{3 \times 3}$ are the
proportional, integral, and derivative gain matrices of the position loop.
The scalar thrust is obtained by projecting $\mathbf{F}_i^{\mathrm{des}}$ onto
the current body $z$-axis, and the desired attitude
$\boldsymbol{\Phi}_i^{\mathrm{des}}$ is extracted from the direction
of~$\mathbf{F}_i^{\mathrm{des}}$.

\paragraph{Attitude loop.}
A second PID controller produces the desired body-frame torque:
\begin{equation}
  \boldsymbol{\tau}_i^{\mathrm{des}} =
    \mathbf{K}_p^{\mathrm{att}}\,\mathbf{e}_{\Phi}
    + \mathbf{K}_i^{\mathrm{att}}\!\int\!\mathbf{e}_{\Phi}\,\mathrm{d}t
    + \mathbf{K}_d^{\mathrm{att}}\,\mathbf{e}_{\omega},
  \label{eq:pid_att}
\end{equation}
where $\mathbf{e}_{\Phi}$ is the rotation error between the current attitude
$\boldsymbol{\Phi}_i$ and the desired attitude
$\boldsymbol{\Phi}_i^{\mathrm{des}}$,
$\mathbf{e}_{\omega}$ is the angular rate error, and
$\mathbf{K}_p^{\mathrm{att}}, \mathbf{K}_i^{\mathrm{att}},
\mathbf{K}_d^{\mathrm{att}} \in \mathbb{R}^{3 \times 3}$ are the attitude PID
gain matrices.

\paragraph{Motor mixer.}
The scalar thrust and desired torques are combined through a mixer matrix
$\mathbf{M} \in \mathbb{R}^{4\times3}$ to obtain the four motor commands:
\begin{equation}
  \boldsymbol{\omega}^{\mathrm{rpm}} =
    \alpha_{\mathrm{pwm}}\,\mathrm{clip}\!\bigl(
      T_{\mathrm{base}}\,\mathbf{1}_4
      + \mathbf{M}\,\boldsymbol{\tau}_i^{\mathrm{des}},\;
      \omega_{\min},\, \omega_{\max}
    \bigr)
    + \beta_{\mathrm{pwm}},
  \label{eq:mixer}
\end{equation}
where $T_{\mathrm{base}}$ is the hover thrust command,
$\mathbf{1}_4 = (1,1,1,1)^{\top}$ distributes it equally to four motors,
$\mathrm{clip}(\cdot,\, \omega_{\min},\, \omega_{\max})$ clamps the PWM signal
within the hardware limits $[\omega_{\min},\, \omega_{\max}]$, and
$\alpha_{\mathrm{pwm}}, \beta_{\mathrm{pwm}}$ are the affine coefficients of
the PWM-to-RPM conversion.

\subsubsection{Observation and Coverage Indicators}

Each UAV monitors entities within its sensing range.
Let $\mathbf{p}_j^{(t)}$ denote the position of entity~$E_j$ at time step~$t$.
The observation indicator function of~$U_i$ is defined as
\begin{equation}
  I_{\mathrm{obs}}(U_i, E_j, t) =
  \begin{cases}
    1 & \text{if } \lVert\mathbf{p}_i^{(t)} - \mathbf{p}_j^{(t)}\rVert
        \le R_{\mathrm{obs}}, \\
    0 & \text{otherwise},
  \end{cases}
  \label{eq:obs_indicator}
\end{equation}
where $I_{\mathrm{obs}}(U_i, E_j, t) = 1$ indicates that $U_i$ detects
entity~$E_j$ at time~$t$.
Let $\mathbf{p}_m^{(t)}$ denote the position of target~$G_m$.
The coverage indicator function is
\begin{equation}
  I_{\mathrm{cov}}(U_i, G_m, t) =
  \begin{cases}
    1 & \text{if } \lVert\mathbf{p}_i^{(t)} - \mathbf{p}_m^{(t)}\rVert
        \le R_{\mathrm{cov}}, \\
    0 & \text{otherwise},
  \end{cases}
  \label{eq:cov_indicator}
\end{equation}
where $I_{\mathrm{cov}}(U_i, G_m, t) = 1$ indicates that target~$G_m$ is
successfully covered by~$U_i$.
The objective is to maximise the cumulative coverage over the episode
horizon~$T$:
\begin{equation}
  J = \sum_{i=1}^{N}\sum_{t=0}^{T} I_{\mathrm{cov}}(U_i, G_m, t).
  \label{eq:coverage_obj}
\end{equation}
Upon coverage, the target is immediately respawned at a random collision-free
location to keep the task complexity constant.

\subsubsection{Dynamic Entities}

Each target $G_m$ is modelled as a cylinder of radius~$r_{\mathrm{tgt}}$
located at a fixed altitude~$z_{\mathrm{tgt}}$.
The target moves on the $xy$-plane at a constant speed~$v_{\mathrm{tgt}}$;
its heading is re-sampled uniformly from $\mathcal{U}(0, 2\pi)$ every
$T_{\mathrm{hold}}$~control steps.
Each obstacle zone $T_k$ is a cylinder of radius
$r_k \sim \mathcal{U}(r_{\min}^{\mathrm{obs}}, r_{\max}^{\mathrm{obs}})$,
where $r_{\min}^{\mathrm{obs}}$ and $r_{\max}^{\mathrm{obs}}$ are the minimum
and maximum allowable obstacle radii.
Obstacle zones move on the $xy$-plane at speed~$v_{\mathrm{obs}}$ following
the same random-heading protocol as targets.
Boundary reflections keep all entities within the task region.

\subsubsection{Safety Constraints}

To avoid collisions, UAVs must maintain safe distances.
Let $\mathbf{p}_k^{(t)}$ denote the centre of obstacle zone~$T_k$ at
time~$t$.
Let $R_{\mathrm{safe}}^{1} = 2\,r_{\mathrm{uav}}$ denote the minimum safe
distance between UAVs, and
$R_{\mathrm{safe}}^{2} = r_{\mathrm{uav}} + r_k$ the minimum safe distance
between a UAV and an obstacle zone boundary.
The constraints are
\begin{equation}
  \begin{cases}
    \lVert\mathbf{p}_i^{(t)} - \mathbf{p}_j^{(t)}\rVert
      \ge R_{\mathrm{safe}}^{1},
    & \forall\, U_i, U_j \in \mathcal{U},\; i \ne j, \\[4pt]
    \lVert\mathbf{p}_i^{(t)} - \mathbf{p}_k^{(t)}\rVert
      \ge R_{\mathrm{safe}}^{2},
    & \forall\, U_i \in \mathcal{U},\;
      \forall\, T_k \in \mathcal{T}.
  \end{cases}
  \label{eq:safety}
\end{equation}
Each UAV is additionally surrounded by a threat zone of
width~$d_{\mathrm{threat}}$; when another entity enters this zone, a proximity
penalty is incurred (see Section~\ref{subsubsec:reward}).

% ==============================================================
\subsection{Problem Transformation}
\label{subsec:problem_transformation}

The cooperative multi-UAV target coverage task is modelled as a
\emph{Decentralised Partially Observable Markov Decision Process}~(Dec-POMDP),
defined by the tuple
$\langle N, \mathcal{S}, \mathcal{O}, \mathcal{A}, \mathcal{P}, R, \gamma
\rangle$,
where $N$~is the number of UAVs,
$\mathcal{S}$ is the global state space,
$\mathcal{O}$ is the observation space,
$\mathcal{A}$ is the action space,
$\mathcal{P}(\mathbf{s}' \mid \mathbf{s}, \mathbf{a})$ is the state transition
function,
$R: \mathcal{S} \times \mathcal{A} \to \mathbb{R}$ is the reward function, and
$\gamma \in (0,1]$ is the discount factor.
The global state is $\mathbf{s} \in \mathcal{S} = \mathbb{R}^{N \times D}$
with state dimension~$D$.
Each UAV~$U_i$ receives a local observation
$o_i = O(\mathbf{s}) \in \mathbb{R}^{d}$ with $d < D$,
$a_i \in \mathcal{A}$ is its action and
$\mathbf{a} = (a_1, \dots, a_N)$ the joint action.
Each UAV follows a decentralised policy~$\pi_{\theta_i}(a_i \mid o_i)$ with
learnable parameters~$\theta_i$, aiming to maximise the expected discounted
return $G_t = \sum_{k=0}^{\infty} \gamma^k r_{t+k}$, where $r_{t+k}$ is the
reward received at time $t+k$.

\subsubsection{Observation Space}
\label{subsubsec:observation}

At each time step~$t$, $U_i$ receives a local observation $o_i^{(t)}$ composed
of four parts:
\begin{enumerate}
  \item \emph{Self-state}
    $[\mathbf{p}_i^{(t)},\, \mathbf{v}_i^{(t)}]$:
    the position and velocity of~$U_i$, normalised by the region half-extent
    and maximum speed, respectively.
  \item \emph{Teammate state}
    $[\mathbf{p}_j^{(t)} - \mathbf{p}_i^{(t)},\,
      \mathbf{v}_j^{(t)} - \mathbf{v}_i^{(t)}]$,
    $j = 1,\dots,N,\, j \ne i$:
    the relative position and velocity of other UAVs with respect to~$U_i$.
  \item \emph{Target state}
    $[\mathbf{p}_m^{(t)} - \mathbf{p}_i^{(t)}]$,
    $m = 1,\dots,M$:
    the relative position of each target with respect to~$U_i$.
  \item \emph{Obstacle zone state}
    $[\mathbf{p}_k^{(t)} - \mathbf{p}_i^{(t)},\, r_k]$,
    $k = 1,\dots,K$:
    the relative position and radius of each obstacle zone.
    The radius is normalised by $R_{\mathrm{obs}}$.
\end{enumerate}
For entities beyond the observation radius~$R_{\mathrm{obs}}$, the
corresponding entries are set to the zero vector, indicating no observation.

\subsubsection{Action Space}
\label{subsubsec:action}

The action of~$U_i$ is a two-dimensional velocity direction on the $xy$-plane,
$a_i^{(t)} = (a_i^{x}, a_i^{y}) \in [-1,1]^{2}$.
The policy~$\pi_{\theta_i}(a_i \mid o_i)$ outputs normalised actions, which are
scaled by $v_{\max}$ and converted to rotor RPMs by the cascaded PID controller
described in Section~\ref{subsec:system_model}.

\subsubsection{Reward Function}
\label{subsubsec:reward}

The reward function decomposes into two components: collision penalty and
coverage reward.

\paragraph{Collision penalty.}
Each UAV is surrounded by a threat zone of width~$d_{\mathrm{threat}}$.
For a UAV pair $(U_i, U_j)$ with distance
$d = \lVert\mathbf{p}_i - \mathbf{p}_j\rVert$:
\begin{equation}
  r_1 =
  \begin{cases}
    -\alpha_{\mathrm{col}} & 0 \le d \le 2\,r_{\mathrm{uav}}, \\[3pt]
    \displaystyle -\frac{\beta}{\exp(d - 2\,r_{\mathrm{uav}})}
    & 2\,r_{\mathrm{uav}} < d < 2\,r_{\mathrm{uav}} + d_{\mathrm{threat}},
    \\[3pt]
    0 & \text{otherwise},
  \end{cases}
  \label{eq:r1}
\end{equation}
where $\alpha_{\mathrm{col}}$ is the hard-collision penalty coefficient applied
when two UAVs physically collide, and $\beta$ is the soft threat-zone
penalty coefficient that decays exponentially with distance.

Similarly, for UAV~$U_i$ and obstacle zone~$T_k$ with planar distance
$d = \lVert\mathbf{p}_{i,xy} - \mathbf{p}_{k,xy}\rVert$, where
$\mathbf{p}_{i,xy}$ and $\mathbf{p}_{k,xy}$ are the $xy$-projections of the
UAV and obstacle zone positions:
\begin{equation}
  r_2 =
  \begin{cases}
    -\alpha_{\mathrm{col}} & 0 \le d \le D_k, \\[3pt]
    \displaystyle -\frac{\beta}{\exp(d - D_k)}
    & D_k < d < D_k + d_{\mathrm{threat}}, \\[3pt]
    0 & \text{otherwise},
  \end{cases}
  \label{eq:r2}
\end{equation}
where $D_k = r_{\mathrm{uav}} + r_k$ is the contact distance between the UAV
body and the obstacle zone boundary.

\paragraph{Coverage reward.}
Each time a UAV covers a target (i.e.\
$I_{\mathrm{cov}}(U_i, G_m, t) = 1$), the nearest covering UAV receives a
positive reward $r_3 = \alpha_{\mathrm{cov}}$, where $\alpha_{\mathrm{cov}}$ is
the coverage reward coefficient.

The total reward for each UAV is
\begin{equation}
  R_{\mathrm{UAV}} = r_1 + r_2 + r_3.
  \label{eq:total_reward}
\end{equation}

\subsubsection{Episode Termination}
\label{subsubsec:termination}

An episode terminates when:
(i)~the cumulative number of captures reaches
$\lceil \eta \cdot M \rceil$, where $\eta \in (0,1]$ is the coverage
completion ratio;
(ii)~any safety constraint in~\eqref{eq:safety} is violated; or
(iii)~the episode duration exceeds the maximum time horizon~$T_{\max}$.
